generating the canonical filtration.

Writing @@M0@@, with @@M1@@ and the @@M2@@ monomials in the @@M3@@, and using the additivity of @@M4@@, we see that it suffices to give the proof for @@M5@@, with @@M6@@ and @@M7@@.

Using Lemma 4.11. and the @@M8@@-@@M9@@-homomorphism from @@M10@@ to @@M11@@ which sends @@M12@@ to @@M13@@, we see that it suffices to give the proof for an element of the form @@M14@@.

Finally, since the polynomial @@M15@@ has no isobaric components of weight @@M16@@, we are reduced to the case @@M17@@, which itself follows from the fact that the polynomials @@M18@@ have no isobaric components of weight @@M19@@ (3.6).

\sgaStableTarget{sga6.e05.rem.4-13}\textbf{Remarks 4.13.} i) The decomposition of @@M20@@ given in 4.12 is a direct-sum decomposition of @@M21@@-modules. In other words, if @@M22@@, and if we write @@M23@@, with @@M24@@, @@M25@@, and @@M26@@, then @@M27@@.

ii) Statement 4.12. remains true for every quotient of @@M28@@ by a @@M29@@-ideal contained in the augmentation ideal.

iii) The statement analogous to 4.12., obtained by replacing the indeterminate @@M30@@ by a family of indeterminates @@M31@@, is also true.

\sgaStableTarget{sga6.e05.cor.4-14}\textbf{Corollary 4.14.} \sgaStableTarget{sga6.e05.eq.4-14-1}@@M32@@.

We define an arrow from right to left by associating to each element @@M33@@, where @@M34@@ is the class in @@M35@@ of an element @@M36@@ of @@M37@@ and @@M38@@, the class in @@M39@@ of @@M40@@. We define an arrow in the opposite direction by decomposition according to Proposition 4.12 and Remark 4.13.i). It is clear that these two arrows are inverse to one another.

